\appendix

\chapter{Master arm firmware}
\label{app:firmware}

The complete sketch is at \repo{firmware/master\_arm/teensy\_final.ino}. The
superseded revision, which carried the pin conflict described in
\cref{ch:master}, is kept at
\repo{firmware/master\_arm/teensy\_merged\_prior.ino} so the change can be
diffed.

\noindent\emph{The inertial-sensor driver header is included by the sketch but
was not among the files handed over, and is not in the repository. It is
therefore not listed below. Every call site that uses it appears in the
listings that follow, so the interface the system depends on is fully
specified here even though the implementation is not.}

\section{Pin assignment as declared}

\begin{lstlisting}[language=C++, caption={Pin arrays. \texttt{A4} and
\texttt{A5} appear in neither array, which is the fix described in
\cref{ch:master}.}]
const int potPinsK1[JOINTS] =
{
 A0, A1, A2, A3, A6, A7, A8      // K1 (left arm) -- 7 channels, no I2C pins
};
const int potPinsK2[JOINTS] =
{
 A9, A10, A11, A12, A13, A16, A17  // K2 (right arm) -- 7 channels, no I2C pins
};

const int FSR1_PIN = A14;
const int FSR2_PIN = A15;

const int BUTTON1_PIN = 2;
const int BUTTON2_PIN = 3;
\end{lstlisting}

\section{Signal conditioning as implemented}

\begin{lstlisting}[language=C++, caption={The potentiometer filter and the
count-to-degree transform. The filter operates in raw ADC units; note that
\texttt{setup()} seeds the accumulator with \texttt{analogRead(pin) << 8},
which is inconsistent with this recurrence and produces the start-up transient
quantified in \cref{tab:transient}.}]
static constexpr int ALPHA_FP    = 64;
static constexpr int ADC_MIN_VAL = 50;
static constexpr int ADC_MAX_VAL = 4045;

inline int updateEMA(int pin, int32_t& emaFP)
{
    int raw = analogRead(pin);
    emaFP =
    ((int32_t)ALPHA_FP * raw +
    (256 - ALPHA_FP) * emaFP) >> 8;
    return (int)emaFP;
}

inline float adcToDeg(int adc)
{
    float t = constrain(
    (float)(adc - ADC_MIN_VAL) /
    (float)(ADC_MAX_VAL - ADC_MIN_VAL),
    0.0f,
    1.0f);
    return t * DEG_MAX;
}
\end{lstlisting}

\begin{lstlisting}[language=C++, caption={Button debounce and toggle. The
output is a latched state, so a host that drops frames cannot miss a press.}]
static constexpr uint32_t DEBOUNCE_MS = 40;

inline int updateButtonToggle(
int pin, int& state, int& lastRaw, uint32_t& lastChangeMs)
{
    int raw = digitalRead(pin);
    uint32_t now = millis();
    if(raw != lastRaw && (now - lastChangeMs) > DEBOUNCE_MS)
    {
        lastChangeMs = now;
        lastRaw = raw;
        if(raw == LOW)
        {
            state = state ? 0 : 1;
        }
    }
    return state;
}
\end{lstlisting}

\section{Frame construction}

\begin{lstlisting}[language=C++, caption={Frame assembly. Inertial values are
transmitted raw in LSB, so the scale factors are applied on the host and can
be corrected without reflashing.}]
    for(int j=0;j<JOINTS;j++)
        p += snprintf(p, end-p, "k1j%d:%.1f,", j+1, degK1[j]);
    for(int j=0;j<JOINTS;j++)
        p += snprintf(p, end-p, "k2j%d:%.1f,", j+1, degK2[j]);
    p += snprintf(p, end-p, "K1IMU:%d,%d,%d,%d,%d,%d,", ax1,ay1,az1,gx1,gy1,gz1);
    p += snprintf(p, end-p, "K2IMU:%d,%d,%d,%d,%d,%d,", ax2,ay2,az2,gx2,gy2,gz2);
    p += snprintf(p, end-p, "fsr1:%d,fsr2:%d,", fsr1, fsr2);
    p += snprintf(p, end-p, "btn1:%d,btn2:%d", btn1, btn2);
\end{lstlisting}

\chapter{Supporting records}
\label{app:records}

The following are maintained in the repository rather than reproduced here,
because each is longer than the chapter that cites it and each is generated
rather than written.

\begin{table}[htbp]
  \centering
  \caption{Supporting records and where they live.}
  \label{tab:records}
  \begin{tabular}{@{}p{52mm}p{90mm}@{}}
    \toprule
    Record & Content \\
    \midrule
    \repo{thesis\_report/DIVERGENCE.md} &
      Divergence from the planning report, by objective, with the reason and
      the evidence for each. \\
    \repo{thesis\_report/EVIDENCE.md} &
      Every numerical claim in this thesis with the file and run that produced
      it. \\
    \repo{recordings/baselines/} &
      Channel-health baselines by date, workspace sweeps, and the five-task
      verification result. \\
    \repo{docs/NEXT\_SESSION.md} &
      Current blocking items in priority order. \\
    \bottomrule
  \end{tabular}
\end{table}

\chapter{Evaluation criteria}
\label{app:criteria}

Each criterion carries a threshold and the justification for that threshold. A
criterion with a number and no justification for the number is not a
criterion. Two criteria that a reader might expect are deliberately absent and
the reason is given at the end.

\begin{table}[H]
  \centering
  \footnotesize
  \caption[Evaluation criteria]{The nine criteria this work is judged against,
  with the outcome of each.}
  \label{tab:criteria}
  \begin{tabular}{@{}p{5mm}p{28mm}p{58mm}p{38mm}@{}}
    \toprule
    & Criterion & Threshold, and its justification & Outcome \\
    \midrule
    1 & Convention round trip &
      Converting a joint angle from the manufacturer's convention to the
      middleware's and back must return the original to within numerical
      precision. The failure mode is a near-full revolution of a joint beside
      a person, so there is no tolerance to spend &
      \textbf{Met} \\
    2 & Directional fidelity &
      Under \SI{30}{\degree} per axis. Beyond that an operator must
      consciously correct rather than move naturally &
      \textbf{Partly met}: vertical and fore-and-aft pass, lateral does not
      (\cref{ch:results}) \\
    3 & Command tracking in simulation &
      Under \SI{5}{\milli\metre} RMS. Below that the residual is dominated by
      the pose deadband rather than by the controller &
      \textbf{Met}, \SIrange{0.7}{4.7}{\milli\metre} \\
    4 & Wearer clearance &
      At least the configured floor at every point of every path, measured
      geometrically &
      \textbf{Not met} on three of four task paths; reported per task in
      \cref{fig:clearance} \\
    5 & Safety response completeness &
      All injected faults detected, responded to and made externally visible.
      A mechanism handling thirteen of fourteen is not a safety mechanism &
      \textbf{Met}, 14 of 14 \\
    6 & Degraded-mode integrity &
      Under \SI{1}{\milli\metre} of commanded motion contributed by failed
      channels from a still hand, and a refusal when no positional channel
      survives &
      \textbf{Partly met}: right arm \SI{0.014}{\milli\metre}, left arm
      \SI{54.5}{\milli\metre}; refusal implemented \\
    7 & Scenario reachability &
      Ten repeats of ten over the whole densified path. A pose feasible on
      \SIrange{72}{82}{\percent} of calls passes a five-of-five check
      \SI{19}{\percent} of the time, so anything short of $N/N$ is not an
      established pose &
      \textbf{Met}, all verified scenarios \\
    8 & Control-path latency &
      Under \SI{200}{\milli\second}, inherited from the planning report &
      \textbf{Not measured end to end}; the control-path figures in
      \cref{ch:results} are lower bounds \\
    9 & Object handling verified in the recording &
      The fingers must open, close within the holding band and reopen, and the
      object must track the gripper &
      \textbf{Met}, all recorded grasp cells \\
    \bottomrule
  \end{tabular}
\end{table}

\textbf{No force or impedance criterion.} The mock hardware returns
identically zero for every force and effort field, so any quantity derived
from a force reading is \emph{structurally undefined} on this platform rather
than merely unmeasured. Force reflection and closed-loop overshoot cannot be
defined as criteria until either a dynamic simulator or the physical arm is in
the loop.

\textbf{No angular-tracking criterion.} The commanded orientation is pinned to
the workspace anchor and never follows the master's wrist, because the wrist
channels are dead or incoherent. There is no angular tracking to measure until
those channels are repaired.

\chapter{Equipment}
\label{app:equipment}

\begin{table}[H]
  \centering
  \small
  \caption[Equipment]{Hardware and software, with the reference for each.}
  \label{tab:equipment}
  \begin{tabular}{@{}p{34mm}p{100mm}@{}}
    \toprule
    Item & Detail \\
    \midrule
    Manipulators & Two Kinova Gen3, seven degrees of freedom,
      \SI{8.2}{\kilo\gram} each, \SI{0.902}{\metre}
      reach~\cite{kinova2022gen3} \\
    Grippers & Robotiq two-finger, \SI{85}{\milli\metre} stroke \\
    Wrist cameras & Colour and depth, one per arm \\
    Frame & Backpack harness; arms mounted behind the shoulders with forward
      tilt \\
    Master arm & Scale model of the same joint sequence: fourteen rotary
      potentiometers, two inertial units, two force-sensitive resistors, two
      buttons \\
    Microcontroller & Teensy 4.1~\cite{pjrc2020teensy41}, \SI{100}{\hertz}
      frame rate \\
    Inertial sensors & MPU-6050 at power-on
      defaults~\cite{invensense2013mpu6050,invensense2013mpuregmap} \\
    Controllers & Two tracked hand controllers, headset stationary as the
      tracking reference (\cref{sec:desk}) \\
    Host & Laptop, Linux, \SI{4}{\gibi\byte} of video memory \\
    Middleware & ROS~2~\cite{macenski2022ros2} \\
    Inverse kinematics & TRAC-IK~\cite{beeson2015tracik} \\
    Motion generation & Ruckig~\cite{berscheid2021ruckig} \\
    Path planning & RRT-Connect~\cite{kuffner2000rrtconnect} in joint space \\
    Plane fitting & RANSAC~\cite{fischler1981ransac}, constrained to the
      support normal \\
    Detection & YOLO-World~\cite{cheng2024yoloworld} with
      FastSAM~\cite{zhao2023fastsam} \\
    Body tracking & MediaPipe Pose~\cite{lugaresi2019mediapipe} \\
    Prior simulation & MuJoCo~\cite{todorov2012mujoco}; Isaac
      Lab~\cite{isaaclab} evaluated and rejected \\
    \bottomrule
  \end{tabular}
\end{table}

\chapter{Where the evidence lives}
\label{app:evidence}

Every quantity quoted in this thesis is produced by a named program and stored
as a record, so that it can be recomputed rather than taken on trust. The
mapping below is the index; the records themselves are in the repository.

\begin{table}[H]
  \centering
  \small
  \caption[Evidence index]{Principal claims and the programs that produce
  them.}
  \label{tab:evidence}
  \begin{tabular}{@{}p{58mm}p{76mm}@{}}
    \toprule
    Claim & Produced by \\
    \midrule
    Channel health verdicts &
      \repo{trajectory\_capture/channel\_report.py} \\
    Cost of each rung of the degradation ladder &
      \repo{scripts/measure\_capability\_ladder.py} \\
    Commanded motion injected by failed channels &
      \repo{scripts/prove\_degraded\_mode.py} \\
    Continuous reachability by direction &
      \repo{scripts/measure\_workspace.py} \\
    Both-arm overlap and the dead band &
      \repo{scripts/sweep\_mount\_overlap.py} \\
    Mount parameter sweep &
      \repo{scripts/sweep\_mount\_geometry.py} \\
    The centre, asked four ways &
      \repo{scripts/search\_centre\_on\_surface.py},
      \repo{measure\_centre\_vs\_height.py},
      \repo{measure\_centre\_vs\_wearer\_posture.py} \\
    Clearance along each task path &
      \repo{scripts/measure\_clearance\_region.py} \\
    Positioning error budget &
      \repo{scripts/measure\_control\_budget.py} \\
    Motion generation and the step sweep &
      \repo{scripts/measure\_teleop\_motion.py} \\
    Simulation-to-hardware terminal error &
      \repo{scripts/measure\_sim\_to\_real\_gap.py} \\
    Map accuracy &
      \repo{scripts/calibrate\_environment.py} \\
    Wearer tracking rate and fallbacks &
      \repo{scripts/measure\_wearer\_tracking\_rate.py},
      \repo{verify\_wearer\_fallbacks.py} \\
    Instruction handling &
      \repo{scripts/sweep\_t1\_instructions.py} \\
    Fault injection &
      \repo{srl\_teleop/fault\_injector.py} \\
    Operating window audit &
      \repo{scripts/verify\_gui\_buttons.py} \\
    Test totals &
      \repo{scripts/check\_tests.py} \\
    \bottomrule
  \end{tabular}
\end{table}

\chapter{Participant documentation}
\label{app:documents}

\noindent\fbox{\parbox{0.95\linewidth}{\small\itshape
To be attached before submission: the participant information sheet, both
consent forms, the pre-session screening questions and the debrief script. Two
consent forms rather than one, because the two roles carry materially
different risks: one participant carries \SI{17}{\kilo\gram} and is within
reach of a manipulator moving under another participant's command, and the
other is not. Drafts exist in the project's research records and are not
reproduced here because they have not been through ethical review.}}
